\begin{abstract}
Extended reality (XR) puts the network inside a perceptual loop. When the user's head turns, the new view
must appear within a few tens of milliseconds. Once content, computation or other participants are
remote, the network must carry 360-degree and volumetric media at high rates, return rendered frames
inside that deadline, take a continuous uplink of pose and sensor data, and keep several people and
machines consistent. This survey organises networked XR by the six functions that every such system
performs. It must \emph{encode} the scene into something that can be sent, \emph{anticipate} what the
user will need next, decide where to \emph{compute}, \emph{provision} the network, \emph{connect} people
and machines that share a world, and \emph{assure} that the experience is good enough and safe. For each
function we state the design question and group the approaches into families by their underlying idea.
We then separate what has been established on real systems from what is contested and what is merely
assumed. Two couplings tie the functions together: a shared deadline that every function spends from,
and a trade between network and device that every design shifts one way or the other. Across all six
functions the same weaknesses recur. Objectives are rarely validated against what users perceive,
requirement figures are inherited rather than measured, and controllers are seldom tested beyond the
simulator. The uplink, multi-user scale and the latency cost of security are studied least. We close
with nine cross-cutting gaps and a dependency-ordered research agenda.
\end{abstract}
